\section{The second quadratic--Tate residue class at 269}
\label{sec:cubic-second-residue-strict-vanishing}

The same original-field method gives strict vanishing on the other
admissible quadratic--Tate residue class.  Its class input is a new
certified quotient bound for a degree-six field.  Only that quotient
bound is used; a full class-group equality and a fundamental-unit
basis are unnecessary.

Keep $p=269$, $K=\mathbb Q(\sqrt{-3})$, $\psi=\chi_{-4}$,
the two cubic labels $j=1,2$, the actual Tate character $c_p$,
and the exact strict local conditions of (G114.2).
Every coefficient DVR below is the original valuation ring $O$ of
an arbitrary finite extension $E_p/\mathbb Q_p$ containing the cubic
values, with uniformizer $\varpi$ and residue field $k$.
For an integer $e\ge1$ put
\[
       a_j=\chi^{-j}\psi^{-1}c_p^{-e}|_{G_K},
       D_j=(E_p/O)(a_j).
 \tag{G115.1}
\]

\begin{theorem}[AC-G115-residual: the second original residual group is zero]
\label{thm:ACG115residual}
If $e\equiv133\pmod{268}$, then for both cubic labels and every
finite residue field containing their values,
\[
              \operatorname{Sel}_{\rm str}(K,k(\bar a_j))=0.
 \tag{G115.2}
\]
This is a direct finite arithmetic conclusion.  It has no
automorphic endpoint, auxiliary host, deformation-model, or
Bloch--Kato premise.
\end{theorem}

\begin{proof}
\emph{1. The exact quadratic character and the degree-six field.}
Here $-(e+1)\equiv134\pmod{268}$, so
\[
       \eta_j:=\bar a_j/\bar c_p
          =\bar\chi^{-j}\psi\bar c_p^{134}.
 \tag{G115.3}
\]
The quadratic character $\bar c_p^{134}$ cuts out
$\mathbb Q(\sqrt{269})$.
One can check the sign without a Frobenius convention: if
$\zeta$ is a primitive $p$th root of unity and
$\tau=\sum_{a\bmod p}(a/p)\zeta^a$, then an automorphism
$\zeta\mapsto\zeta^b$ sends $\tau$ to $(b/p)\tau$.
The double-sum identity gives $\tau\bar\tau=p$;
since $p\equiv1\pmod4$, $\bar\tau=\tau$ and $\tau^2=p$.
Indeed in the double sum, setting $a=bt$ makes the inner sum
over $b\ne0$ equal to $p-1$ for $t=1$ and $-1$ otherwise,
which gives $p$ by the sum-zero identity for the quadratic
character.  This proves the stated quadratic field directly.

Thus $\vartheta=\psi\bar c_p^{134}$ is the character of
$\mathbb Q(\sqrt{-269})$.  Over $K$ it equals the restriction
of the real quadratic character of
$R=\mathbb Q(\sqrt{807})$, since multiplication by the
character of $K/\mathbb Q$ changes the radicand to $807$.
Put $t=\sqrt[3]{269}$, $s=\sqrt{807}$, and
\[
       u=s/t,\qquad L=\mathbb Q(t,s)=\mathbb Q(u),
       \qquad M=K(t,s)=K(u).
 \tag{G115.4}
\]
The identities
\[
             u^2=3t,\qquad u^3=3s,\qquad u^6=7263
 \tag{G115.5}
\]
prove both field equalities.  The polynomial
$Y^6-7263$ is Eisenstein at $269$, hence irreducible.
The real quadratic field $R$ is not the unique quadratic
subfield $K$ of the $S_3$ extension $K(t)/\mathbb Q$.
Consequently
\[
 \operatorname{Gal}(M/\mathbb Q)=S_3\times C_2,
 \qquad \Delta=\operatorname{Gal}(M/K)=C_3\times C_2,
 \tag{G115.6}
\]
and $M$ is exactly the field cut out by $\eta_j$.
The element $b$ of $C_2$ fixes $K(t)$ and sends $s$ to $-s$,
so it sends $u$ to $-u$; the target Kummer character
$\eta_j^{-1}$ is nontrivial cubic with $b$-sign $-1$.
This is the $D=807$ field of (G103.8), not the $D=3$ field.

Over $K_p$, the same polynomial is Eisenstein because $K_p$
is the unramified quadratic extension of $\mathbb Q_p$.
It follows that $M$ has a unique $p$-place, with ramification
index six and residue degree two over $\mathbb Q_p$.
All of $\Delta$ is inertia there.
Normalize its valuation by $v(u)=1$, hence $v(p)=6$.

\emph{2. The certified class quotient and the target class part.}
The new PARI 2.19.1 input is exactly $Y^6-7263$.
Its retained output has signature $(2,2)$, discriminant
$2433913492059072$, and computed cyclic class group $C_{14}$.
The separate, successful call \texttt{bnfcertify(b,1)} ends with
\[
             \texttt{CLASS\_QUOTIENT\_CERTIFY=1},
                       \qquad\texttt{DONE}.
 \tag{G115.7}
\]
The literal official manual conclusion for this flag is that
the actual class group is a quotient of the computed group.
It follows unconditionally that
\[
                         \operatorname{Cl}(L)[269]=0.
 \tag{G115.8}
\]
No equality with $C_{14}$ and no completeness of the displayed
units is claimed.  A separate default-flag certification was
intentionally stopped after (G115.7) succeeded; its partial log
and the termination receipt are preserved and are not treated
as either a success or a failed certificate.

Let $r$ be the reflection of the $S_3$ factor fixing $t$.
It fixes $s$ as well, so $M^{\langle r\rangle}=L$.
Degree-two extension and norm give
\[
       (\operatorname{Cl}(M)[p]\otimes k)^{\langle r\rangle}=0
 \tag{G115.9}
\]
by (G115.8).  A nonzero vector $x$ in the target
$\eta_j^{-1}$ eigenspace would give a nonzero invariant
$x+rx$, since $rx$ belongs to the other cubic eigenspace.
This contradicts (G115.9), so the target class part is zero.
The same holds for the $S$-class group, by applying the integral
prime-to-$p$ character projector to the finite $p$-primary class
group and its quotient; here $S$ consists of the unique $p$-prime.

The logarithmic unit theorem and prime-to-$p$ idempotents give
the regular $\Delta$ representation minus the trivial character
on units modulo $p$, exactly as in G114.
There is no $p$-torsion unit because $[M:\mathbb Q]=12<268$.
The $S$-unit prime-permutation term is trivial, and therefore
does not change either target cubic character.
Thus the $b=-1$, nontrivial-cubic unit part has dimension two
over $\mathbb F_{269}$ and is simple as an
$\mathbb F_{269}[C_3]$ module.  Each of its two character
lines over $k$ has dimension one.

\emph{3. An explicit nonzero local unit image.}
In $L=\mathbb Q(u)$ take the element
\[
\begin{split}
 U={}&466357309795346854156
       +\frac{317922111743885429405}{3}u\\
   &+24081304710076734947u^2
       +\frac{49249635842423393791}{9}u^3\\
   &+\frac{3730460073071408017}{3}u^4
       +\frac{7629317537790237701}{27}u^5.
\end{split}
 \tag{G115.10}
\]
Exact multiplication in the basis $1,u,\ldots,u^5$, using
$u^6=7263$, gives its monic integral characteristic polynomial
\[
 Z^6-AZ^5+BZ^4-CZ^3+BZ^2-AZ+1,
 \tag{G115.11}
\]
where
\[
\begin{split}
 A&=2798143858772081124936,\\
 B&=2272692936963916701196079597387551887,\\
 C&=811803899632075359719052256808584701208252890193904.
\end{split}
\]
The retained exact-arithmetic script constructs that matrix
and substitutes its polynomial in the six-dimensional algebra.
The result is zero.  Monicity and constant term one prove that
$U$ and $U^{-1}$ are integral, so $U$ is an actual global unit
without assuming anything about a computed unit basis.

Put
\[
                   v=U/b(U),\qquad
                   w=v^3/N_{M/K(s)}(v).
 \tag{G115.12}
\]
Then $b(w)=w^{-1}$ and $N_{C_3}(w)=1$.
Thus the class of $w$ lies in the simple unit module of Step~2.
Since $b(u)=-u$ and the coefficients of $U$ are rational,
reduction at the $p$-place gives
\[
 U=A_0+B_0u\pmod{u^2},\quad
 A_0=268,\quad B_0=98,\quad B_0/A_0=171
                               \quad\text{in }\mathbb F_p.
 \tag{G115.13}
\]
Hence $v$ has constant term one and linear coefficient
$2\cdot171$.  The cubic norm has constant term one and no
linear term: the sum of the three cubic roots of unity is zero.
It follows that
\[
                    w=1+219u\pmod{u^2},
                  \qquad219\equiv6\cdot171\pmod{269}.
 \tag{G115.14}
\]
This coefficient is nonzero.

Here $v(p)=6<p-1$, so logarithm and exponential identify
$1+\mathfrak m$ with $\mathfrak m$, and multiplication by
$p$ has image $\mathfrak m^7$.
Since the residue-unit group has order $p^2-1$ prime to $p$,
(G115.14) is not a local $p$th power.
The equivariant localization of the simple two-dimensional
unit module is therefore nonzero, hence injective.
It is injective on both target character lines after extension
of scalars to $k$.

Finally, the same exact Kummer descent as (G114.7) identifies
the strict group with the $\eta_j^{-1}$ part of
\[
 \ker\left(V_{M,S}\otimes k\longrightarrow
                      \bigoplus_{w\mid p}M_w^\times/M_w^{\times p}
                                      \otimes k\right).
 \tag{G115.15}
\]
The target class part is zero by Step~2, and the unit localization
is injective by Step~3.  This proves (G115.2), with full local
zero at $p$ and every away-inertia condition retained.
\end{proof}

\begin{theorem}[AC-G115-strict: both admissible quadratic--Tate classes]
\label{thm:ACG115strict}
For $\psi=\chi_{-4}$, every integer $e\ge1$ with
\[
                         e\equiv133\text{ or }267\pmod{268},
 \tag{G115.16}
\]
both $j=1,2$, and every original finite coefficient DVR $O$
containing the cubic values,
\[
                \operatorname{Sel}_{\rm str}(K,(E_p/O)(a_j))=0.
 \tag{G115.17}
\]
In particular its length and normalized length over that same
$O$ are zero.  No ramification-index-one hypothesis is required.
\end{theorem}

\begin{proof}
For $e\equiv267$, this is AC-G114-strict.
For $e\equiv133$, use AC-G115-residual and the direct
uniformizer argument of G114, whose local checks do not depend
on which of the two residue classes is used.
More explicitly, at $p$ raising an inertia element with
nontrivial cubic value to the $268$th power kills its
rational residual cyclotomic value but preserves its cubic
value.  Since $\psi$ is unramified there,
$D_j^{G_{K_p}}=D_j^{G_K}=0$.
At each $v\nmid p$, $a_j|_{I_v}$ has order dividing six,
so its invariant module is zero or all of the divisible
module $D_j$.
The sequence
$0\to k(\bar a_j)\to D_j\xrightarrow{\varpi}D_j\to0$
therefore lifts every strict $\varpi$-torsion class to a
strict residual class, which is zero.
Every continuous cocycle $G_K\to D_j$ has finite image
because its source is compact and its target discrete, so
each class is killed by a power of $\varpi$.
A nonzero strict class would have a nonzero $\varpi$-torsion
multiple.  This contradiction proves (G115.17) over the
original $O$, without any coefficient extension.
\end{proof}

\paragraph{Certificate and geometric scope.}
The new class input is the successful quotient-only
\texttt{bnfcertify(b,1)} run for $Y^6-7263$, with input,
output and binary data retained.  Its conclusion is exactly
that the class group is a quotient of $C_{14}$.
The separate exact unit script proves the unit and its local
image independently of any claimed fundamental-unit basis.
The separate default certification was stopped intentionally
once these sufficient inputs were available.

The pure Kummer, unit-representation and coefficient-kernel
arguments have the same previously checked source scope as
G103 and G114.  Geometrically, (G115.15) is the specified
character part of Kummer torsors on
$\operatorname{Spec}\mathcal O_M[1/p]$ trivial over the
$p$-adic completion.  The local unit calculation and class
quotient make that part zero, and the original-uniformizer
argument passes to the discrete coefficient system.
No automorphic endpoint, auxiliary host, model M4, mixed
characteristic-zero Bloch--Kato vanishing, other coefficient
prime, other Dirichlet character or relaxed local condition
is supplied or assumed by this result.
